\documentclass[10pt,a4paper]{amsart}
\usepackage{amsmath,amssymb,amsthm,graphicx}
% Retained preamble for the Stacks Project chapter morphisms.tex (Morphisms of Schemes).
% The chapter file pinned-original-source.tex carries only \input{preamble} on its first line, so no
% preamble text lives inside the chapter itself.  The Stacks Project uses one single shared preamble for
% every chapter; the snapshot stages it verbatim in sources/stacks-chow/preamble.tex.  The block below is
% copied mechanically, line for line, from that file's lines 14-18 (the Xy-pic setup) and lines 154-217
% (the native theorem declarations together with the mathematical macros).  Nothing was added, removed or
% redefined.  The wrapper class and the amsmath/amssymb/amsthm/graphicx base packages are supplied by the
% assembler, so the class/\externaldocument part of the shared preamble is not repeated here.
% For commutative diagrams we use Xy-pic
\usepackage[all]{xy}

% We use 2cell for 2-commutative diagrams.
\xyoption{2cell}
% Theorem environments.
%
\theoremstyle{plain}
\newtheorem{theorem}[subsection]{Theorem}
\newtheorem{proposition}[subsection]{Proposition}
\newtheorem{lemma}[subsection]{Lemma}

\theoremstyle{definition}
\newtheorem{definition}[subsection]{Definition}
\newtheorem{example}[subsection]{Example}
\newtheorem{exercise}[subsection]{Exercise}
\newtheorem{situation}[subsection]{Situation}

\theoremstyle{remark}
\newtheorem{remark}[subsection]{Remark}
\newtheorem{remarks}[subsection]{Remarks}

\numberwithin{equation}{subsection}

% Macros
%
\def\lim{\mathop{\mathrm{lim}}\nolimits}
\def\colim{\mathop{\mathrm{colim}}\nolimits}
\def\Spec{\mathop{\mathrm{Spec}}}
\def\Hom{\mathop{\mathrm{Hom}}\nolimits}
\def\Ext{\mathop{\mathrm{Ext}}\nolimits}
\def\SheafHom{\mathop{\mathcal{H}\!\mathit{om}}\nolimits}
\def\SheafExt{\mathop{\mathcal{E}\!\mathit{xt}}\nolimits}
\def\Sch{\mathit{Sch}}
\def\Mor{\mathop{\mathrm{Mor}}\nolimits}
\def\Ob{\mathop{\mathrm{Ob}}\nolimits}
\def\Sh{\mathop{\mathit{Sh}}\nolimits}
\def\NL{\mathop{N\!L}\nolimits}
\def\CH{\mathop{\mathrm{CH}}\nolimits}
\def\proetale{{pro\text{-}\acute{e}tale}}
\def\etale{{\acute{e}tale}}
\def\QCoh{\mathit{QCoh}}
\def\Ker{\mathop{\mathrm{Ker}}}
\def\Im{\mathop{\mathrm{Im}}}
\def\Coker{\mathop{\mathrm{Coker}}}
\def\Coim{\mathop{\mathrm{Coim}}}

% Boxtimes
%
\DeclareMathSymbol{\boxtimes}{\mathbin}{AMSa}{"02}

%
% Macros for moduli stacks/spaces
%
\def\QCohstack{\mathcal{QC}\!\mathit{oh}}
\def\Cohstack{\mathcal{C}\!\mathit{oh}}
\def\Spacesstack{\mathcal{S}\!\mathit{paces}}
\def\Quotfunctor{\mathrm{Quot}}
\def\Hilbfunctor{\mathrm{Hilb}}
\def\Curvesstack{\mathcal{C}\!\mathit{urves}}
\def\Polarizedstack{\mathcal{P}\!\mathit{olarized}}
\def\Complexesstack{\mathcal{C}\!\mathit{omplexes}}
% \Pic is the operator that assigns to X its picard group, usage \Pic(X)
% \Picardstack_{X/B} denotes the Picard stack of X over B
% \Picardfunctor_{X/B} denotes the Picard functor of X over B
\def\Pic{\mathop{\mathrm{Pic}}\nolimits}
\def\Picardstack{\mathcal{P}\!\mathit{ic}}
\def\Picardfunctor{\mathrm{Pic}}
\def\Deformationcategory{\mathcal{D}\!\mathit{ef}}

% --- appended 2026-10-12 (Chief): the project own definitions of the reference/slogan/history
% display environments, copied mechanically from the shared Stacks preamble lines 8-12
% (\usepackage{verbatim} plus the three \newenvironment lines, which the chapter preamble keeps
% verbatim from the project preamble). Nothing added, removed or redefined: these environments
% carry no mathematical content and the project itself typesets them as comments. Without these
% three lines a counted statement whose own text uses them fails with Environment slogan undefined.
\usepackage{verbatim}
\newenvironment{reference}{\comment}{\endcomment}
\newenvironment{slogan}{\comment}{\endcomment}
\newenvironment{history}{\comment}{\endcomment}
\usepackage{cleveref}
\crefname{originallemma}{original lemma}{original lemmas}
\Crefname{originallemma}{Original lemma}{Original lemmas}
\makeatletter
\newlabel{morphisms-lemma-finite-type-over-affine-ample-very-ample}{{01VT}{source}{}{sourcepointer-1}{}}
\newlabel{morphisms-lemma-finite-type-over-affine-ample-very-ample@cref}{{[originallemma][0][]01VT}{[1][0][]source}{}{}{}}
\makeatother
\crefname{originallemma}{original lemma}{original lemmas}
\Crefname{originallemma}{Original lemma}{Original lemmas}
\makeatletter
\newlabel{lemma-ample-over-affine}{{01VK}{source}{}{sourcepointer-2}{}}
\newlabel{lemma-ample-over-affine@cref}{{[originallemma][0][]01VK}{[1][0][]source}{}{}{}}
\makeatother
\crefname{originallemma}{original lemma}{original lemmas}
\Crefname{originallemma}{Original lemma}{Original lemmas}
\makeatletter
\newlabel{lemma-quasi-projective-finite-type-over-S}{{01VS}{source}{}{sourcepointer-3}{}}
\newlabel{lemma-quasi-projective-finite-type-over-S@cref}{{[originallemma][0][]01VS}{[1][0][]source}{}{}{}}
\makeatother
\crefname{originallemma}{original lemma}{original lemmas}
\Crefname{originallemma}{Original lemma}{Original lemmas}
\makeatletter
\newlabel{constructions-lemma-projective-space-bundle}{{01OE}{source}{}{sourcepointer-4}{}}
\newlabel{constructions-lemma-projective-space-bundle@cref}{{[originallemma][0][]01OE}{[1][0][]source}{}{}{}}
\makeatother
\crefname{originallemma}{original lemma}{original lemmas}
\Crefname{originallemma}{Original lemma}{Original lemmas}
\makeatletter
\newlabel{lemma-ample-very-ample}{{01VN}{source}{}{sourcepointer-5}{}}
\newlabel{lemma-ample-very-ample@cref}{{[originallemma][0][]01VN}{[1][0][]source}{}{}{}}
\makeatother
\crefname{originallemma}{original lemma}{original lemmas}
\Crefname{originallemma}{Original lemma}{Original lemmas}
\makeatletter
\newlabel{lemma-ample-power-ample}{{02NN}{source}{}{sourcepointer-6}{}}
\newlabel{lemma-ample-power-ample@cref}{{[originallemma][0][]02NN}{[1][0][]source}{}{}{}}
\makeatother
\title{Six equivalent formulations of ampleness over an affine base in finite type}
\author{The Stacks Project Authors}
\date{stacks-project:morphisms.tex:local-snapshot-20261012; GNU FDL 1.2 or later}
\begin{document}
\maketitle
\setcounter{subsection}{1}
\noindent\textit{Editorial note.} All delivered mathematical text is the author's own, copied from the pinned Stacks Project chapter file morphisms.tex; no statement or proof step was rewritten, reordered or removed.\par\medskip
\noindent\textit{Editorial note.} The chapter file pinned-original-source.tex carries only \textbackslash{}input\{preamble\} on its first line, so no preamble text lives inside the chapter. The retained file sources/stacks-morphisms/retained-preamble.tex is a mechanical, line-for-line copy of the project's single shared preamble as staged in the snapshot. Nothing was added, removed or redefined, and the excerpt serves only as preamble, never as a span coordinate.\par\medskip
\noindent\textit{Editorial note.} Both counted spans start at the \textbackslash{}begin line of their environment, so the statement and the author's proof are delivered whole; surrounding prose of the section is deliberately outside the counted spans and is carried as standard prerequisites instead.\par\medskip
\noindent\textit{Editorial note.} The citation fields (source\_id, source\_version, authors, source\_url, license.name, source\_title) are identical to those of every other card this lane has produced from this chapter, so the batch selector treats them as one source.\par\medskip
\noindent\textit{Editorial note.} Read in full under the Lead's 12:4x rule. The proof contains 'Trivially (6) implies (5)' / 'Trivially (4) implies (3)' - the word marks genuinely trivial implications between the listed conditions, and every non-trivial step is a cited result, so the assembly of the equivalence is the author's own work.\par\medskip
\noindent\textit{Allowed standard prerequisites (editorial).}\begin{itemize}
\item Ampleness on an affine total space equals relative ampleness. Lemma \textbackslash{}ref\{lemma-ample-over-affine\}, cited at proof line 9538. Identifies the first two conditions.
\item Relative ampleness gives an immersion into projective space for finite type morphisms over affine bases. Lemma \textbackslash{}ref\{lemma-quasi-projective-finite-type-over-S\}, cited at proof lines 9539-9540. Gives condition (6).
\item Very ampleness implies ampleness, and ampleness is stable under taking powers. Lemma \textbackslash{}ref\{lemma-ample-very-ample\} and Lemma \textbackslash{}ref\{lemma-ample-power-ample\}, cited at proof lines 9548-9549. Close the cycle.
\end{itemize}\medskip
\ifdefined\hypertarget\hypertarget{sourcepointer-1}{}\fi
\noindent\textit{Source reference.} Original lemma 01VT: focused-seeds/AGP/sources/stacks-snapshot-20261011/tags.txt, source lines 2458--2458. The lemma's own canonical tag, quoted so that readers can resolve the label inside the delivered text; it points at the counted statement itself and carries no proof content.\par\medskip
\ifdefined\hypertarget\hypertarget{sourcepointer-2}{}\fi
\noindent\textit{Source reference.} Original lemma 01VK: focused-seeds/AGP/sources/stacks-snapshot-20261011/tags.txt, source lines 2449--2449. Navigation pointer to a statement the counted proof refers to but does not contain; nothing load-bearing is delegated to it.\par\medskip
\ifdefined\hypertarget\hypertarget{sourcepointer-3}{}\fi
\noindent\textit{Source reference.} Original lemma 01VS: focused-seeds/AGP/sources/stacks-snapshot-20261011/tags.txt, source lines 2457--2457. Navigation pointer to a statement the counted proof refers to but does not contain; nothing load-bearing is delegated to it.\par\medskip
\ifdefined\hypertarget\hypertarget{sourcepointer-4}{}\fi
\noindent\textit{Source reference.} Original lemma 01OE: focused-seeds/AGP/sources/stacks-snapshot-20261011/tags.txt, source lines 2191--2191. Navigation pointer to a statement the counted proof refers to but does not contain; nothing load-bearing is delegated to it.\par\medskip
\ifdefined\hypertarget\hypertarget{sourcepointer-5}{}\fi
\noindent\textit{Source reference.} Original lemma 01VN: focused-seeds/AGP/sources/stacks-snapshot-20261011/tags.txt, source lines 2452--2452. Navigation pointer to a statement the counted proof refers to but does not contain; nothing load-bearing is delegated to it.\par\medskip
\ifdefined\hypertarget\hypertarget{sourcepointer-6}{}\fi
\noindent\textit{Source reference.} Original lemma 02NN: focused-seeds/AGP/sources/stacks-snapshot-20261011/tags.txt, source lines 3486--3486. Navigation pointer to a statement the counted proof refers to but does not contain; nothing load-bearing is delegated to it.\par\medskip
\begin{lemma}
\label{lemma-finite-type-over-affine-ample-very-ample}
Let $f : X \to S$ be a morphism of schemes.
Let $\mathcal{L}$ be an invertible $\mathcal{O}_X$-module.
Assume $S$ affine and $f$ of finite type.
The following are equivalent
\begin{enumerate}
\item $\mathcal{L}$ is ample on $X$,
\item $\mathcal{L}$ is $f$-ample,
\item $\mathcal{L}^{\otimes d}$ is $f$-very ample for some $d \geq 1$,
\item $\mathcal{L}^{\otimes d}$ is $f$-very ample for all $d \gg 1$,
\item for some $d \geq 1$ there exist $n \geq 1$ and an immersion
$i : X \to \mathbf{P}^n_S$ such that
$\mathcal{L}^{\otimes d} \cong i^*\mathcal{O}_{\mathbf{P}^n_S}(1)$, and
\item for all $d \gg 1$ there exist $n \geq 1$ and an immersion
$i : X \to \mathbf{P}^n_S$ such that
$\mathcal{L}^{\otimes d} \cong i^*\mathcal{O}_{\mathbf{P}^n_S}(1)$.
\end{enumerate}
\end{lemma}

\begin{proof}
The equivalence of (1) and (2) is Lemma \ref{lemma-ample-over-affine}.
The implication (2) $\Rightarrow$ (6) is
Lemma \ref{lemma-quasi-projective-finite-type-over-S}.
Trivially (6) implies (5).
As $\mathbf{P}^n_S$ is a projective bundle over $S$ (see
Constructions, Lemma \ref{constructions-lemma-projective-space-bundle})
we see that
(5) implies (3) and (6) implies (4) from the definition of a
relatively very ample sheaf.
Trivially (4) implies (3). To finish we have to show that
(3) implies (2) which follows from Lemma \ref{lemma-ample-very-ample}
and Lemma \ref{lemma-ample-power-ample}.
\end{proof}

\end{document}
